\section{How to read the formal companion}
\label{sec:reading}
Chapter V places several kinds of mathematics in one signal processing language: coordinate partitions, bounded Hilbert space operators, finite filter banks, finite tree optimization, formal series, and projective actions. The formal development keeps their types and hypotheses distinct. Following those distinctions is the most useful way to read its proofs.

A pointwise partition identity needs little structure. An energy theorem needs a normed space with the intended norm. A transform theorem needs convergence. An optimal finite pruning needs a finite family and an objective. A fractional linear expression needs a chart, while its projective action survives at the chart's poles. Each transition has a named construction or theorem in the source.

\begin{keyidea}[title={Follow the model before the theorem name}]
Start with the ambient type, then locate the bridge to the intended signal model. A coordinate identity, a bounded operator theorem, and an analytic transform law have different inputs. The exact excerpt shows where those inputs enter the declaration.
\end{keyidea}

\subsection{What the completed contract contains}
The selected release, \code{1.1.0-expanded-lean-checked}, records twenty scoped theorem-like manuscript results with explicit proofs and 225 supporting theorem targets. The inherited namespace audit reports 827 declarations. These are different denominators: the 827 includes definitions and compiler-generated declarations, and the 225 includes supporting results rather than 225 independent new mathematical claims.

The formal core is the expanded manuscript preserved as Part I of the final traditional edition. Part II preserves the broader Annals development. Its additional proof records are prose mathematics. They do not enter the twenty-result formalization contract. Section~\ref{sec:extension} gives a separate crosswalk.

The two edge recovery conjectures remain open. Known-edge cancellation is established, measurable-mask perturbation is established, and finite adaptive selection is established within its hypotheses. Those results do not provide an estimator with the full statistical and approximation guarantees requested by the conjectures.

\subsection{Evidence available in this edition}
\begin{longtable}{>{\raggedright\arraybackslash}p{.24\textwidth}>{\raggedright\arraybackslash}p{.66\textwidth}}
\toprule Evidence & Meaning \\ \midrule\endhead
Source reconciliation & Current hashes, 225 declaration locations, 225 audit directives, and 225 inherited axiom-output records agree. \\ \addlinespace
Historical Lean execution & The September 6 UTC release records the pinned build and namespace audit passing. The raw logs are preserved. \\ \addlinespace
Manual correspondence & Manuscript statements have a reviewed, explicit relationship to formal declarations. Lean does not parse or certify the English PDF automatically. \\ \addlinespace
Exact model tests & Finite independent models exercise indexing, masks, trees, lifting, filters, and boundary cases. They are computational evidence on their tested ranges. \\ \addlinespace
Companion numerics & Signal experiments and figure data are empirical demonstrations. Their floating-point implementation is outside the Lean theorem claim. \\ \bottomrule
\end{longtable}
The guide build performs no Lean compilation. Reproduction commands appear in Section~\ref{sec:reproduce}, with the missing-tool case treated as a failure to establish a fresh kernel result. The final package's verification records distinguish current reruns from inherited evidence.

\subsection{Read a theorem as a typed contract}
The name after \code{theorem} is a reusable declaration. Parenthesized arguments are explicit inputs. Braced arguments are generally inferred. A hypothesis such as \code{(hN : 0 < N)} is an actual input to the theorem. A typeclass argument such as \code{[AddCommMonoid E]} specifies the algebraic laws available to its proof.

The expression after the colon is the proposition. The proof after \code{:= by} constructs evidence of that proposition. Common commands here have small, specific jobs: \code{ext} reduces equality to coordinates, \code{rw} rewrites using an equality, \code{simp} uses declared simplification rules, \code{ring} normalizes polynomial identities, \code{omega} handles integer arithmetic, and \code{linarith} combines linear inequalities. Their successful use belongs to the recorded compilation, not to a textual count of tactic names.

A definition marked \code{noncomputable} may use classical choice or real-number constructions that do not give executable numerical code. Its theorems can still be checked by the kernel. In particular, the pruning evaluator over exact real costs specifies a mathematical algorithm and a count of unit operations. Lean evaluation speed is outside that count.

\subsection{Four kinds of signal space}
\begin{longtable}{>{\raggedright\arraybackslash}p{.25\textwidth}>{\raggedright\arraybackslash}p{.29\textwidth}>{\raggedright\arraybackslash}p{.35\textwidth}}
\toprule Object & Lean representation & What it provides \\ \midrule\endhead
Arbitrary sequence & \code{ℤ → K} & Every integer sample exists. No norm or convergence is assumed. \\ \addlinespace
Finite filter & \code{ℤ →₀ K} & A finitely supported coefficient function. Convolution sums are finite. \\ \addlinespace
Square-summable signal & \code{lp (fun _ : I => ℂ) 2} & The joint $\ell^2$ norm, infinite norm-convergent expansions, bounded maps and adjoints. \\ \addlinespace
Formal series & \code{R⟦X⟧} & Coefficients with formal addition and Cauchy multiplication. Analytic convergence is not part of the type. \\ \bottomrule
\end{longtable}
The same informal word ``packet'' appears in all four settings. The guide always identifies the ambient type before using an identity across them.

\subsection{Recommended reading routes}
For the signal processing route, read integer addresses, the Hilbert bridge, FIR transfer, faithful filter banks, and Haar packets. For the adaptive route, read known edges, mask perturbations, complete prunings, the dynamic program, and its stability envelope. For the algebraic route, read formal cyclic products, homogeneous Smith actions, and CRT masks.

Every displayed code excerpt is copied from an immutable source declaration. The raw excerpt, its source line range, and its SHA256 hash accompany the PDF. Excerpts rely on the imports, namespaces, variables, and definitions in their source files. They are not advertised as independently compilable snippets. The complete project remains the authoritative executable object.
